\section{Residual Class Under Revealed-Sacrifice Observation}
\label{sec:residual_class}

\subsection{Re-characterization}

A natural alternative approach would define the residual class as
capabilities failing benchmarkability (identity-individuated,
non-repeatable, non-communicable).  The revealed-sacrifice framing
produces a differently-shaped residual class that is, by our
interpretation, easier to apply within sacrifice-channel
diagnostics: the channel-reach criterion is directly checkable
on observable trade infrastructure, while the
benchmarkability criterion requires structural classification
of the capability before it is observed.  Whether this
constitutes ``operationally cleaner'' is an interpretation,
not a theorem.

\begin{definition}[Residual Class Under Revealed Sacrifice]
\label{def:residual_class}
The residual class is the set of capabilities \emph{structurally
outside the reach} of the sacrifice channel---capabilities whose
exercise, by their nature, produces no sacrifice event even with
unlimited observation time:
\[
  \ResS = \{c \in P : \text{no sacrifice-observable path exists
  in principle}\}.
\]
$\ResS$ is a \emph{structural} property of the capability, not a
property of the observation window.  A capability is in $\ResS$
because the sacrifice channel \emph{cannot} reach it (purely
private exercise with no material footprint), not because the
channel has not yet observed it.  Establishing $\ResS$
membership is a \emph{structural/domain proof obligation}, not
a directly checkable property of trade infrastructure: the
analyst must show that every conceivable exercise mode of the
candidate capability fails to produce a sacrifice event, which
involves enumeration over hypothetical exercise pathways
rather than inspection of observed trades.  Trade
infrastructure shows what happened; $\ResS$ membership claims
what cannot happen in principle.  Examples are restricted to
capabilities that admit \emph{no} sacrifice-observable exercise
mode in principle: private contemplative experience whose
exercise consumes neither paid material nor time-displaced
opportunity beyond idle cycles, and passively-concurrent
internal capabilities such as ongoing interoceptive
self-awareness exercised concurrently with all other activity
(see proof of
Proposition~\ref{prop:residual_intersection}(c) for a worked
witness).  Cognitive capabilities that admit paid-training,
dedicated-time, or counterparty-supported exercise modes
(silent rehearsal as a paid course, planning sessions on
deliberately blocked time, supervised reasoning practice) lie
\emph{outside} $\ResS$ because at least one
sacrifice-observable mode exists in principle, even if no
specific window observes it.
\end{definition}

\begin{remark}[Structural versus evidential absence]
$\ResS$ is distinct from the set of capabilities with no sacrifice
evidence in a given window.  A capability outside $\ResS$
(tradable in principle) may have no sacrifice evidence in a
particular window simply because no trade happened to cover
it---this is the \emph{dormant} category in the gap decomposition
(Definition~\ref{def:gap_decomp}), not the residual class.
Dormant capabilities can be resolved by waiting for more trade
data; residual capabilities cannot.
\end{remark}

\begin{remark}[Conservatism of $\ResS$ classification]
\label{rem:residual_conservatism}
Classifying a capability as residual is a structural claim that
should be conservative.  The criterion ``any plausible
mechanism'' is intentionally strict, but ``plausible'' is
itself contestable: ``private contemplative experience'' might
displace idle cycles that have time-opportunity cost,
``baseline'' awareness capabilities might be measurable through
clinical assessment that costs the agent's time.  We adopt the
strict reading: classify $c \in \ResS$ only when every
proposed sacrifice path can be shown to belong to a
\emph{distinct capability} (e.g., the act of being assessed is a
distinct capability from the assessed underlying state) under
the framework's poset-construction discipline.  The strict
reading shrinks $\ResS$ relative to a looser reading; this is
the conservative direction (a smaller residual class makes the
gap-decomposition's residual cell smaller, which is
preferable since $\delta_{\mathrm{residual}}$ is the
irreducible component beyond the channel's reach).  In
practice, residual-class membership claims should be
accompanied by an explicit argument that all candidate
sacrifice paths attach to distinct capabilities, not the target
$c$ itself.  When such arguments are absent, classify $c$ as
non-residual by default.

If there is any plausible mechanism by
which exercise of~$c$ could produce a sacrifice event (even
indirectly---e.g., an artistic practice that requires purchasing
supplies), $c$ should be classified as outside $\ResS$.  The
residual class should be small by construction.

\emph{$\ResS$ is an analyst-dependent governance parameter,
not a structural fact certifiable from trade infrastructure.}
The ``in principle'' criterion of
Definition~\ref{def:residual_class} requires a domain analyst
to adjudicate which proposed sacrifice paths attach to the
target capability and which to distinct capabilities; reasonable
analysts will disagree on cases (e.g., does ``baseline
interoceptive self-awareness'' admit a clinical-assessment
sacrifice path that attaches to it, or does the assessment
attach to a distinct ``measured-awareness'' capability under
the individuation discipline?).  Proposition~\ref{prop:residual_intersection}'s
conditional clause makes this explicit: whether
$\ResS \setminus \ResB$ is non-empty depends on individuation
choices.  $\delta_{\mathrm{residual}}$ should therefore be
read as a governance-parametrized quantity whose magnitude
encodes the deployment's residual-class adjudication
discipline, not as a structural fact independent of analyst
choice.  Two deployments operating on the same trade pool
under different individuation conventions can produce
different $\delta_{\mathrm{residual}}$; the gap decomposition
itself is robust to the choice (the additive identity holds
under any consistent classification), but cross-deployment
$\delta_{\mathrm{residual}}$ comparisons require the
classification conventions to be stated.
\end{remark}

\subsection{Relationship to the benchmarkability-based residual class}

\begin{proposition}[Residual Class Intersection]
\label{prop:residual_intersection}
Let $\ResB$ be the benchmarkability-based residual class
(capabilities failing benchmarkability under a direct-measurement
approach) and $\ResS$ be the revealed-sacrifice residual class
(Definition~\ref{def:residual_class}).  Then:
\begin{enumerate}
\item[(a)] $\ResB \setminus \ResS$ is non-empty (witnessed in
  (b) below).
\item[(b)] Capabilities in $\ResB \setminus \ResS$:
  identity-individuated capabilities that nonetheless involve
  observable sacrifice.  Example: an artistic practice that is
  identity-individuated but involves purchasing supplies, paying
  for studio space, foregoing labour hours.  The time-sacrifice
  channel captures a lower bound even though the capability is
  identity-individuated.
\item[(c)] Whether $\ResS \setminus \ResB$ is non-empty is
  \emph{individuation-dependent}: under poset-construction
  choices that distinguish baseline from trained/honed forms
  of internal capabilities (treating them as distinct nodes),
  candidate witnesses such as baseline interoceptive
  self-awareness lie in $\ResS \setminus \ResB$; under
  individuation choices that treat baseline and trained as
  co-extensive, no witness in $\ResS \setminus \ResB$ has
  been established by this paper.  See (c) of the proof.
\item[(d)] $\ResB \cap \ResS$ is non-empty (witnessed in (d)
  below): identity-individuated capabilities exercised purely
  privately with no sacrifice event.  Example: private
  contemplative experience with no material input or
  time-sacrifice signal.
\end{enumerate}
The proposition's downstream uses in this paper require only
(a), (b), and (d); the conditional (c) is presented for
completeness.  The classical and channel-reach residual carves
therefore overlap (witnessed by (d)) but are not in a
containment relationship in the strict direction $\ResB \subseteq
\ResS$ (witnessed by (b)); the converse direction
$\ResS \subseteq \ResB$ is not asserted by this paper.
\end{proposition}

\begin{proof}
Each part is established by exhibiting a capability with the
claimed membership, which proves non-containment in the relevant
direction.

\emph{(a) $\ResB \setminus \ResS$ non-empty (witnessed in (b))
ruling out $\ResB \subseteq \ResS$.}
Part (b) below produces an explicit witness for
$\ResB \setminus \ResS \neq \emptyset$.  The reverse
direction $\ResS \subseteq \ResB$ is not established by this
proposition (whether it holds depends on the conditional
witness construction in (c) below); the proposition therefore
asserts only the asymmetric statement that $\ResB$ is not a
subset of $\ResS$, plus the intersection-non-empty witness in
(d).

\emph{(b) $\ResB \setminus \ResS$.}
Consider the capability ``a specific artist's studio practice.''
Under the benchmarkability criterion of~\citep[%
Definition~2]{lasser2026poset}, this is identity-individuated: the
artefact sequence is non-repeatable and non-communicable in the
benchmark sense, so the capability lies in $\ResB$.  But the
practice consumes materials (paid-for supplies, rent on studio
space) and displaces outside-option labour hours; every session
admits a money- or time-sacrifice event
(\citep[Revealed~Sacrifice]{lasser2026paper8a}) on the underlying bundle.  Whether
that event ultimately emits a $\volRlower$ signal additionally
requires the sacrifice channel's S0--S5 admissibility and
genuine-exchange conditions to hold for the event; what
$\ResS$-membership turns on is the \emph{reachability in
principle} of a sacrifice-observable path, not the
admissibility of any particular event.  Reachability is
satisfied here because each session
\emph{admits} (in the reachability-in-principle sense---not
that S0--S5 are necessarily verified for each event) a
sacrifice event under standard market conditions; therefore
the capability is reachable in principle by the sacrifice
channel and so $\notin \ResS$.

\emph{(c) $\ResS \setminus \ResB$.}
The residual definition (Definition~\ref{def:residual_class})
requires that \emph{no} exercise mode produces a sacrifice event
in principle---not merely that some modes do not.  Many cognitive
capabilities that feel purely internal (silent rehearsal,
planning, reasoning) admit paid-training or dedicated-time
exercise modes and therefore lie outside $\ResS$.  The genuine
witnesses for $\ResS \setminus \ResB$ are capabilities exercised
passively and concurrently with other activity, with no
incremental material input, no incremental time displacement, and
no counterparty.

$\ResS \setminus \ResB$ witnesses are harder to construct
than $\ResB \setminus \ResS$ because clinical measurement,
training programs, and structured-report assessments
typically generate sacrifice events on capabilities that
\emph{seem} purely internal, pulling such capabilities out of
$\ResS$.  We give a candidate witness with explicit
disclosure of the structural argument's load-bearing claim,
which the reader may judge.

Consider ``\emph{baseline} interoceptive self-awareness as a
background condition of conscious embodiment, distinguished
from \emph{trained interoceptive acuity} which clinical
neuroscience can measure.''  Trained acuity is benchmarkable
\emph{and} reachable by the sacrifice channel through
paid-training programs, mindfulness courses with explicit fees,
and clinically-billable assessment sessions, so trained acuity
lies in neither $\ResS$ nor $\ResB$.  The candidate
$\ResS \setminus \ResB$ witness is the \emph{baseline}
component itself.  The load-bearing claim is that
\emph{incrementally exercising the baseline} (as distinct from
training a higher acuity level, or measuring the current
acuity) requires no sacrifice---the baseline is on whenever
the agent is conscious, and any sacrifice associated with
clinical measurement or training accrues to the
\emph{measurement} or \emph{training} capability, not to the
baseline-exercise capability.  This separation is contestable:
a stricter reading would say baseline interoceptive exercise
is co-extensive with trained interoceptive acuity exercise
(any clinical measurement is also a measurement of the
baseline, hence is a sacrifice path to the baseline), in which
case the baseline is reachable in principle and lies outside
$\ResS$.  Whether the separation holds is a domain question
about how interoceptive capability is individuated.  We
adopt the separation: under the framework's poset-construction
discipline, an analyst declaring ``baseline interoceptive
self-awareness'' as a distinct capability node from
``trained interoceptive acuity'' commits to treating
sacrifices for the latter as not-attributable to the former;
this is a poset-construction choice, supported by the
distinction the framework draws between exercise of a
capability and training-level changes to a capability.  Under
this individuation, the baseline witnesses
$\ResS \setminus \ResB$.  An analyst preferring co-extensive
individuation would have to look elsewhere for an
$\ResS \setminus \ResB$ witness; we do not claim the baseline
example is robust against every individuation choice.  Similar
candidate witnesses (subject to the same individuation
assumption) include baseline peripheral sensory integration
and baseline attentional stability.

\emph{(d) $\ResB \cap \ResS$.}
Consider ``private contemplative experience with no material
footprint, exercised during time intervals already
otherwise-priced by the framework's poset construction (so
that no incremental sacrifice attaches to the contemplation
itself).''  This is identity-individuated (non-communicable
phenomenology, so $\in \ResB$) and emits no sacrifice event
under the assumed individuation: no money, no time newly
displaced from sacrifice-attributable activity, no
counterparty, so $\in \ResS$.  The witness depends on the
poset-construction choice that contemplation does not
incrementally displace any sacrifice-attributable activity;
under a stricter reading that attributes opportunity-cost
even to ``idle'' time, the witness fails.  We treat the
witness as conditional in the same sense as
(c): non-empty under the stipulated individuation, and not
asserted as robust against every individuation.  Even with
this conditionality, the proposition's downstream uses
(Section~\ref{sec:gap_decomposition}'s residual cell) only
require $\ResS$ to be \emph{constructible}, not non-empty;
$\ResS = \emptyset$ would just collapse $\delta_{\mathrm{residual}}$
to zero.
\end{proof}

\begin{remark}[The observation-channel carve is the
appropriate one for this paper]
The prior residual class was defined by the \emph{measurement
apparatus} (what benchmarks can test).  The new residual class
is defined by the \emph{observation channel} (what trades
reveal).  The observation-channel carve is the appropriate one
for a paper about what the framework can observe under its
privacy commitments: the question is not ``what could we test
if we had access?'' but ``what do we actually see?''  Other
carves (benchmarkability, computability) remain appropriate
for other questions; we do not claim the
observation-channel carve is correct in the absolute, only
that it is the relevant carve for sacrifice-based diagnostics.
\end{remark}

\subsection{Bounding the residual class}

Paper~3's observational-individuation
result~\citep[Corollary~2.1]{lasser2026horizon} establishes a
\emph{structural} floor on $\volL$: for any distinguishable
agent~$k$, the poset contains an identity capability $d_k$ with
$\OIfloor(k) \geq 1$, so removing~$a_k$ contracts $\volL$ by at
least $\OIfloor(k)$.  The companion paper's $\volRexact(O)$
diagnostic
(the exercise indicator~\citep[Revealed~Sacrifice]{lasser2026paper8a}), by
contrast, sums $\volL$-weights of capabilities $d$ with
$e_t(d) = 1$, where
$e_t(d)$ requires a realized cooperative output in
$[t-\Delta, t]$ that would be unrealizable if $d$ were removed
from $O$.  Transferring the OI floor into $\volRexact(O)$
therefore requires a bridge between the structural $\volL$ fact
and the temporal exercise indicator.  We state that bridge
explicitly as a hypothesis rather than derive it from
Corollary~2.1.

\begin{definition}[OI--Exercise Bridge Condition]
\label{def:oi_exercise_bridge}
An agent~$k$ with OI capability $d_k \in O$ satisfies the
\emph{OI--exercise bridge} at time~$t$ if some realized
cooperative output in $[t-\Delta, t]$ that agent~$k$
participated in is identity-individuating in the sense that
removing $d_k$ from $O$ would make that output unrealizable as
an output attributable to~$k$.  Equivalently,
$e_t(d_k) = 1$ under
the exercise indicator~\citep[Revealed~Sacrifice]{lasser2026paper8a}.  Active presence within
$O$ that is merely dormant (e.g., $k$ is nominally in the
perimeter but contributed nothing to realized cooperative
outputs during $[t-\Delta, t]$) does not satisfy the bridge.
\end{definition}

\begin{proposition}[Active-Agent Floor Transfer, conditional on
the bridge]
\label{prop:oi_floor}
Fix a window $W = [t_0, t_0 + T]$.  For an agent~$k$ with OI
capability $d_k \in O$ satisfying the OI--exercise bridge
(Definition~\ref{def:oi_exercise_bridge}) at some $t \in W$,
$d_k$ is exercised within $W$ and therefore contributes its
$\volL$ weight $w(d_k)$ to $\volRexact^{[W]}(O)$.  Assume
additionally that $w(d_k) \geq \OIfloor(k)$, where $\OIfloor(k)$
is the elimination floor
of~\citep[Corollary~2.1]{lasser2026horizon}.  This weight
premise holds under the binary-default poset construction of
Corollary~2.1 ($w(d_k) = \OIfloor(k) = 1$); under non-binary
weight schemes it must be checked explicitly because
Corollary~2.1 establishes a contraction floor on capability
removal, not a direct lower bound on $w(d_k)$.  Summing over
agents $k \in K$ whose OI capabilities $\{d_k\}$ are pairwise
poset-independent (stated as an explicit hypothesis on the
agent set; the category partition of~\citep[Revealed~Sacrifice]{lasser2026paper8a}
covers the unbenchmarked space $U$, not necessarily the
observation perimeter $O$, so this independence is imposed
directly rather than imported) and who each satisfy the bridge
at some $t \in W$:
\[
  \volRexact^{[W]}(O) \;\geq\; \sum_{k \in K} w(d_k)
  \;\geq\; \sum_{k \in K} \OIfloor(k).
\]
The floor transfers only under both the bridge hypothesis and
the weight premise $w(d_k) \geq \OIfloor(k)$; it is \emph{not}
a consequence of Corollary~2.1 alone, which establishes a
$\volL$ elimination floor rather than an exercise witness or a
direct weight lower bound.
\end{proposition}

\begin{proof}
By the bridge (Definition~\ref{def:oi_exercise_bridge}) at some
$t \in W$, $e_t(d_k) = 1$, so $d_k \in P^{\mathrm{ex}}_t$ for
that $t$.  Under the window-active form of the companion
paper's $\volR$ definition~\citep[Revealed~Sacrifice]{lasser2026paper8a}, $d_k \in
P^{\mathrm{ex}}_W = \bigcup_{s \in W} P^{\mathrm{ex}}_s$, hence
$d_k$ contributes its $\volL$ weight $w(d_k)$ to
$\volRexact^{[W]}(O) = \volL(O \cap P^{\mathrm{ex}}_W)$ (the
window-active sum is $\volL$ on the union of instantaneous
exercised sub-posets, so an exercised capability contributes
its full $w$).  By the weight premise of the proposition,
$w(d_k) \geq \OIfloor(k)$; combining,
$d_k$ contributes at least $\OIfloor(k)$ to
$\volRexact^{[W]}(O)$.  By
Paper~3~\citep[Corollary~2.1]{lasser2026horizon},
$\OIfloor(k) \geq 1$ (the binary default, strengthened by any
subsumption weight in $\volL$).  Pairwise poset-independence of
the $\{d_k\}_{k \in K}$ in the parent poset $P$ lets $\volL$
decompose additively by
axiom~M4~\citep[Proposition~1]{lasser2026poset}, so the
$\volL$ contributions sum without overlap, giving
$\volRexact^{[W]}(O) \geq \sum_{k \in K} w(d_k) \geq
\sum_{k \in K} \OIfloor(k)$.  Without the bridge hypothesis,
Corollary~2.1~\citep{lasser2026horizon} establishes only the
structural $\volL$-on-elimination fact, which does \emph{not}
supply the exercise witness required by the exercise indicator
of the companion paper~\citep[Revealed~Sacrifice]{lasser2026paper8a}.  Without the
weight premise $w(d_k) \geq \OIfloor(k)$, the proof reaches
$\volRexact^{[W]}(O) \geq \sum_{k} w(d_k)$ from the bridge
alone, but the further inequality
$\sum_k w(d_k) \geq \sum_k \OIfloor(k)$ is left unestablished.
\end{proof}

\begin{remark}[When the bridge holds and when it fails]
\label{rem:oi_bridge_scope}
The bridge holds when active participation is
\emph{identity-coupled}: the agent's actions produce realized
cooperative outputs whose attribution depends on the agent
remaining in the perimeter.  Typical operational examples are
participation in a committed trade event, execution of a
verified action with agent-attributable signing, or
contribution to a cooperative capability whose output is signed
by the agent's identifier.  The bridge fails for \emph{dormant
presence}: an agent is nominally within $O$ but contributed
nothing attributable to them in $[t-\Delta, t]$.  In such
windows the structural $\volL$ floor of Corollary~2.1 still
holds (removing the agent would still contract $\volL$), but
$\volRexact(O)$ receives no contribution from $d_k$ until the
agent resumes identity-coupled activity.  This is consistent
with the paper's stance that $\volR$ measures what the
collective \emph{does}, not merely what it possesses.
\end{remark}

\begin{remark}[Channel separation]
Under the boundary-closed hypothesis on $O$
(the exercise indicator~\citep[Revealed~Sacrifice]{lasser2026paper8a}, conditions BC1/BC2),
the OI floor contributes exclusively to $\volRexact(O)$ in the
combined diagnostic (the exercise-indicator appendix of~\citep[Revealed~Sacrifice]{lasser2026paper8a}):
\[
  \volRcomb = \volRexact(O) + \volRlower(P \setminus O).
\]
It does \emph{not} contribute to $\volRlower$ (the
sacrifice-based bound outside $O$).  Conflating the two channels
would double-count.  When boundary-closure fails (either BC1 or
BC2), $\volRcomb$ remains a conservative lower bound on
$\volR(P)$ but undercounts the cross-boundary cooperatives and
subsumption mass
(see the boundary-cooperative-undercount remark
in~\citep[Revealed~Sacrifice]{lasser2026paper8a}); the OI floor's contribution is
unaffected because the OI capability $d_k$ is, by construction,
an internal capability of agent $k$'s substrate---not a
cooperative or subsumption edge spanning the perimeter boundary.
\end{remark}
